% 本文件由 LaTeX 排版 Agent 根据有效 translation.json 自动生成。
% author=Rufinus; work_id=2709; title=护教书

\chapter{护教书}

% structure: p0001 source=h1 target=omitted reason=page-title-already-rendered-by-parent

\Emph{呈递罗马城主教\Person{阿纳斯塔修斯}}\par

本书是因\Person{鲁菲努斯}在就任后不久受到指控而引发的，这些指控是针对\Person{阿纳斯塔修斯}提出的，\Person{阿纳斯塔修斯}自498年至503年担任罗马教座。当时罗马教宗的权威尚未如日后那般强大，因此，\Person{阿纳斯塔修斯}不太可能如某些人所猜测的那样，从\Place{阿奎莱亚}召\Person{鲁菲努斯}前来罗马，以回应正式指控或接受他的审判，因为\Person{鲁菲努斯}当时正与\Person{克罗玛提乌斯}主教以信任关系生活在那里。然而，由于罗马是信息的中心，一位基督徒不会希望被其主教视为不良。那些指控\Person{鲁菲努斯}的人是\Person{耶柔米}在罗马的朋友，尤其是贵族寡妇\Person{马尔塞拉}和元老\Person{帕玛丘斯}。他们曾试图在398年11月\Person{西里修斯}去世前获得对他的一些谴责；但\Person{西里修斯}善待了\Person{鲁菲努斯}（据\Person{耶柔米}说，\Quote{他的纯朴被利用了}）。然而，在399年\Person{阿纳斯塔修斯}当选后，他们指控\Person{鲁菲努斯}因翻译\Person{奥力振}的\WorkTitle{\OriginalTerm{论原理}{\GreekText{Περὶ} \GreekText{᾿Αρχῶν}}}而将异端引入罗马教会。\Person{耶柔米}在\WorkTitle{书信}127:10中这样谈论\Person{马尔塞拉}：\Quote{她是异端被定罪的原因：她带来了从前曾受他们教导、因而沾染错误和异端的证人；她指出有多少人曾被欺骗；她让人把\WorkTitle{\OriginalTerm{论原理}{\GreekText{Περὶ} \GreekText{᾿Αρχῶν}}}的卷册拿来，并指出\OriginalTerm{蝎子}{Scorpion}在其中所作的篡改：直到最后，信件被写出，不止一次，传唤异端者前来辩护，但他们不敢前来。他们所面对的定罪力量如此之大，以致于为了避免异端在他们面前被揭露，他们宁愿缺席而被定罪。}从\Person{阿纳斯塔修斯}致耶路撒冷\Person{若望}关于\Person{鲁菲努斯}的信函中，我们得知，虽然他强烈反对\Person{奥力振}的翻译，但他将\Person{鲁菲努斯}本人留给自己的良心，并不过问他的下落。\Person{鲁菲努斯}的这封信，虽称为\WorkTitle{护教书}，却毫无回应教宗传唤或判决的痕迹，而仅仅是对可能使教宗对他产生偏见的言论的答复。撰写此\WorkTitle{护教书}的年代是公元400年。\par

1. 我得知，某些人在贵圣座管辖区域内就信仰或其他问题挑起争论时，提及了我的名字。我斗胆相信，自幼在教会严格规范中受教的贵圣座，定会拒绝听取任何针对一位缺席者——且是与您在信德和天主之爱中合一、为您所熟知的善意之人——的诽谤。然而，既然我听说有人攻击我的名誉，我认为有必要以书面形式向贵圣座阐明我的立场。我无法亲自前来，因为我离家近三十年后刚刚返回，若如此快便离开那些我久别重逢的亲人，未免过于严苛，几乎不近人情。长途跋涉的辛劳也使我身体虚弱，无力再次启程。我写此信的目的并非为了消除您心中可能存在的疑影——我将您的心灵视为圣地，犹如不容任何邪恶进入的神圣之所。相反，我希望我将要向您作出的宣认，如同一根手杖置于您手中，用以驱散那些如同犬吠般嫉妒我的人。\par

2. 事实上，当异端者迫害我时，我的信德已得到充分证明。当时我寄居在亚历山大教会，因忠心而遭受监禁和流放的刑罚；然而，若有人欲考验我的信德，或愿听闻并了解它，我将予以宣认。我相信天主圣三是一性体、一天主性，具有同一的德能与实体；因此，在圣父、圣子和圣神之间毫无差异，唯独圣父是第一位，圣子是第二位，圣神是第三位。天主圣三是真实且生活的位格，是性体与实体的合一。\par

3. 我也宣认，天主子在末世由童贞女和圣神降生；祂取了我们自然的人性血肉与灵魂；以此受难、被埋葬，并从死者中复活；祂复活时所有的血肉，正是曾安放在坟墓中的同一血肉；并以此同一血肉连同灵魂，在复活后升了天堂，我们期待祂从那里降来审判生者死者。\par

4. 此外，关于我们自身血肉的复活，我相信它将完整而完美；正是我们现在所生活的这同一血肉。我们并不像某些人所诽谤的那样，认为将有另一血肉代替这血肉复活；而是这同一血肉，没有失去任何一个肢体，也没有截去身体的任何一部分；除了可朽性之外，其一切属性无一缺失。这正是圣宗徒关于身体所应许的：\BibleQuote{它播种在可朽坏中，复活在不朽坏中；播种在软弱中，复活在德能中；播种在羞辱中，复活在光荣中；播种的是属生灵的身体，复活的是属神的身体。}这是我从在\Place{阿奎莱亚}教会领受圣洗的那些人那里领受的教义；我认为这也是宗座长期以来所传授和教导的。\par

5. 我肯定，此外，一个审判将要来临，在审判中，每个人将根据他所行的善或恶，接受他肉身生命应得的报应。如果对人的赏报是按照他们的行为，那么对于作为罪之普遍原因的魔鬼，更是如此。关于魔鬼本身，我们的信仰是\BibleBook{}{福音}所记载的，即他和他的所有天使都将承受永火作为他们的份，与他一同承受的还有那些做他的工的人，即那些控告自己弟兄的人。如果有人否认魔鬼将受永火的惩罚，愿他与他一同在永火中有份，这样他就能亲身经历他现在所否认的事实。\par

6. 我接下来被告知，关于灵魂的本性问题引起了一些骚动。这类投诉是否应该被受理还是被搁置，你必须自己决定。然而，如果你希望知道我对这个问题的看法，我将坦率地陈述。我读过许多关于这个问题的作者，我发现他们表达了各种观点。我所读的一些人认为，灵魂是通过人类种子的渠道与物质身体一同注入的；他们给出了一些证明。我认为这是拉丁作者中\Person{戴尔都良}或\Person{拉克坦提乌斯}的观点，也许还有少数其他人。另一些人断言，天主每天都在创造新灵魂，并将它们注入在母腹中形成的身体；而还有一些人相信，灵魂是在很久以前当天主从无中创造万物时就全部造好了，现在他所做的只是按他的意愿将每个灵魂移栽到它的身体中。这是\Person{奥利振}和一些其他希腊人的观点。我自己，在天主面前宣告，在阅读了每一种观点之后，直到目前我无法将其中任何一种视为确定和绝对的；这个问题的真理判定我留给天主，并留给任何他愿意启示的人。因此，我在这点上的宣认是：第一，这些不同的观点是我在书中找到的；但第二，我目前对这个主题仍然无知，除了教会作为信德条款而传递的，即天主是灵魂和身体的创造者。\par

7. 至于另一件事。我被告知有人对我提出反对，因为我应一些弟兄的请求，将\Person{奥利振}的某些著作从希腊文翻译成了拉丁文。我想每个人都明白，这只是出于恶意才将此作为指责的理由。因为，如果作者中有任何冒犯性的陈述，为什么要将其转嫁为译者的过错？我被要求将希腊文中所写的呈现为拉丁文；我所做的只是将拉丁词汇与希腊思想相匹配。因此，如果这些思想中有任何值得称赞的，称赞不属于我；同样，如果有什么应该责备的，责备也不属于我。我承认我在作品中加入了一些我自己的东西；正如我在前言中所说，我运用自己的判断力删除了不少段落；但只删除了那些我怀疑不是\Person{奥利振}本人如此陈述的段落；在我看来，这些陈述是他人插入的，因为我在其他地方发现作者以公教的意义陈述了此事。因此，我恳求你，神圣、可敬且圣善的父，不要因此允许对我兴起恶意的风暴，也不要认可使用党派和诽谤——这些武器在天主的教会中永远不应使用。如果简单的信德和无辜不得到教会的保护，它们在哪里能够安全？我不是\Person{奥利振}的辩护者或捍卫者；我也不是第一个翻译他著作的人。我之前的其他人做了完全同样的事情，我是许多人中最后一位，应弟兄们的请求而做。如果下令不得进行此类翻译，这样的命令适用于未来，而不适用于过去；但是如果那些在发布任何此类命令之前进行翻译的人应该受到责备，那么责备必须从那些迈出第一步的人开始。\par

8. 至于我，我以基督的名宣告，我从未持有，也永远不会持有任何其他信仰，只有上面所陈述的信仰，即由\Place{罗马}教会、\Place{亚历山大}教会以及我自己所在的\Place{阿奎莱亚}教会所持有的信仰；并且\Place{耶路撒冷}也在宣讲这信仰；如果有人相信其他，无论他是谁，愿他受诅咒。但那些仅仅出于恶意和怨恨在弟兄中制造纷争和绊脚石，并使他们跌倒的人，将在审判之日为此交账。\par

\SourceRecord{Translated by W.H. Fremantle.；Nicene and Post-Nicene Fathers, Second Series；Vol. 3.；Edited by Philip Schaff and Henry Wace.；Buffalo, NY: Christian Literature Publishing Co.,；1892.；<http://www.newadvent.org/fathers/2709.htm>.}
